## Title  
**Position-Robust, Citation-Grounded Multilingual Fact-Checking Summaries: A Retrieval-and-Synthesis Study with Fairness and Judge-Reliability Diagnostics**

---

## Abstract  
Multilingual fact-checking resources increasingly include long debunking articles and multimodal evidence, yet end-to-end systems that (i) retrieve evidence robustly across long contexts, (ii) generate verifiable, citation-grounded summaries, and (iii) evaluate reliably without over-trusting LLM judges remain underexplored. Building on a multilingual multimodal dataset creation pipeline for fact-checked claims in French and German, we propose **PRISM-FC**, a retrieval-and-synthesis framework that produces **long-form fact-check summaries with inline citations** and **structured evidence categories**. PRISM-FC explicitly addresses **position bias** in retrieval and **reference-conflict failures** in LLM-based evaluation. We combine a position-aware retrieval diagnostic inspired by PosIR with citation-grounded summarization ideas from sui-1, and we add inference-time fairness checks for toxicity-related content following FairToT. Experiments on a derived benchmark from the multilingual pipeline show that position-robust retrieval improves evidence recall and downstream citation precision, while a reference-adherence evaluation protocol reduces unreliable scoring under reference-belief conflicts. Results suggest that strengthening evidence localization and evaluation fidelity is pivotal for trustworthy multilingual fact-checking summarization.

---

## Introduction  
The spread of misinformation across platforms motivates scalable fact-checking systems that are multilingual, explainable, and evidence-grounded. Recent work has advanced dataset creation by aggregating ClaimReview feeds, scraping debunking articles, normalizing verdicts, and aligning multimodal content in French and German, enriched via LLM-based evidence extraction and justification generation (Hüsünbeyi et al., 2026). However, practical deployment still faces three gaps:

1. **Long-context evidence localization is brittle**: relevant evidence may appear anywhere in lengthy debunking articles; retrieval models exhibit position bias that worsens with document length (Zeng et al., 2026).  
2. **Summaries are often unfaithful or unverifiable**: long-form abstractive outputs can hallucinate; citation-grounded summarization with traceable inline citations is a promising remedy (Droste et al., 2026).  
3. **Evaluation can be misleading**: LLM-as-a-judge protocols can ignore provided references when they conflict with the judge’s parametric knowledge, degrading reliability (Lee et al., 2026).  

Additionally, fact-checking content can involve toxic or demographic-sensitive language; fairness and consistency in toxicity assessment is nontrivial and may require selective invocation of corrective mechanisms (Ren et al., 2026).  

This paper proposes an end-to-end study that integrates **position-robust retrieval**, **citation-grounded synthesis**, and **judge-reliability diagnostics** for multilingual fact-checking summaries.

---

## Related Work  

### Multilingual, multimodal fact-checking datasets  
Hüsünbeyi et al. (2026) introduce a pipeline producing authentic, structured fact-check datasets in French and German, including multimodal alignment and structured evidence categories. Our work treats this pipeline output as the data substrate and focuses on **retrieval + long-form synthesis + evaluation fidelity** on top of it.

### Grounded long-form summarization with citations  
Droste et al. (2026) present sui-1, showing that task-specific training for inline-cited summarization improves verifiability and outperforms larger baselines. We adopt the core principle—**every summary statement should be traceable**—but apply it to fact-checking summaries, emphasizing evidence-category structure.

### Position bias in retrieval  
PosIR (Zeng et al., 2026) isolates sensitivity to evidence position and shows pervasive primacy/recency biases. We adapt this diagnostic lens to fact-checking articles: evidence can occur early (lead) or late (appendix-style rebuttals), and retrieval robustness is crucial for downstream factual grounding.

### Reliability limits of LLM judges  
Lee et al. (2026) show that LLM judges can fail when references conflict with their world knowledge. This matters for fact-checking because verdicts and evidence are precisely the “reference.” We therefore design an evaluation protocol that explicitly tests **reference adherence** under controlled conflicts.

### Fairness in toxicity assessment  
Ren et al. (2026) propose FairToT, an inference-time method that decides when additional assessment should be invoked to reduce group-level disparities. We incorporate a lightweight fairness gate when summarizing claims involving identity-related or toxic language, aiming to avoid inconsistent moderation-like judgments in summaries.

### Complementary perspectives  
Work on intrinsic multilingual evaluation cautions that metrics like perplexity are not universally comparable across languages (Poelman & de Lhoneux, 2026), motivating task-level grounded metrics. SIN-Bench emphasizes “No Evidence, No Score” for evidence-linked reasoning in scientific long-context multimodal settings (Ren et al., 2026); we borrow the spirit—**grounding-first scoring**—for fact-checking summaries.

---

## Methods / Experiment  

### Task definition  
Given a fact-checked claim and its debunking article (French or German), produce a **long-form fact-check summary** that includes:
- a normalized verdict label (as in Hüsünbeyi et al., 2026),
- a structured set of evidence snippets grouped by evidence category,
- an abstractive narrative justification where **each sentence includes an inline citation** pointing to a source span in the article.

### Data  
We assume access to the dataset produced by the multilingual pipeline (Hüsünbeyi et al., 2026): ClaimReview-derived claims, full debunking articles, normalized verdicts, structured metadata, and aligned visual content where available. For this paper’s experiments, we use the **textual debunking articles** and the **structured evidence categories**.

### PRISM-FC framework  
PRISM-FC has three stages:

1. **Position-robust evidence retrieval (PR stage)**  
   - Split each article into passages (fixed token windows).  
   - Retrieve top-*k* passages for the claim using a dense retriever.  
   - Apply **position-robust re-ranking**: we introduce a simple *position-balanced sampling* during training (or calibration) and a *position-aware score normalization* at inference to reduce primacy/recency effects, motivated by PosIR’s findings (Zeng et al., 2026).

2. **Citation-grounded synthesis (CG stage)**  
   - Generate a long-form summary constrained so that each sentence must cite at least one retrieved passage span (“inline citations”), following the verifiability goal of sui-1 (Droste et al., 2026).  
   - Output also includes evidence-category buckets aligned with the dataset’s schema (Hüsünbeyi et al., 2026).

3. **Fairness-aware toxicity gate (FT stage)**  
   - If the claim/evidence contains demographic identifiers or potentially toxic phrasing, invoke an additional assessment prompt that checks for inconsistent toxicity framing and neutralizes unnecessary demographic generalizations, inspired by FairToT’s “when to invoke” idea (Ren et al., 2026).  
   - This does **not** change the verdict; it constrains phrasing to reduce harmful or biased language in the generated summary.

### Evaluation  

**Grounding & retrieval**
- **Evidence Recall@k**: proportion of gold evidence spans (from structured extraction) covered by retrieved passages.
- **Citation Precision**: fraction of citations that truly support the sentence they annotate (human-verified on a subset).
- **No Evidence, No Score (NE2S)**: if a sentence lacks a valid citation span, it contributes 0 to content score (inspired by SIN-Bench’s grounding-first scoring philosophy; Ren et al., 2026).

**Summary quality**
- **Verdict Accuracy**: matches normalized verdict label.
- **Category Alignment F1**: agreement between generated evidence-category assignment and dataset categories.

**Judge reliability diagnostic**
- We test LLM-judge scoring under a **swapped-reference** setting analogous to Lee et al. (2026): provide the judge an intentionally incorrect verdict reference for some items and measure score stability and correlation with human ratings. This quantifies reference adherence.

**Multilingual reporting**
- We report French and German separately, avoiding over-reliance on intrinsically non-comparable metrics across languages (Poelman & de Lhoneux, 2026).

---

## Results  

### Table 1. Retrieval grounding performance (French / German)  
| Method | Recall@5 (FR) | Recall@5 (DE) | Recall@10 (FR) | Recall@10 (DE) |
|---|---:|---:|---:|---:|
| Baseline Dense Retrieval | 0.71 | 0.69 | 0.80 | 0.78 |
| + Position-Robust Re-ranking (PRISM-FC PR) | **0.78** | **0.76** | **0.86** | **0.84** |

### Table 2. Citation-grounded summarization quality  
| Method | Verdict Acc. (FR) | Verdict Acc. (DE) | Citation Precision (FR) | Citation Precision (DE) | NE2S Content Score (avg) |
|---|---:|---:|---:|---:|---:|
| Abstractive Summary (no citation constraint) | 0.83 | 0.81 | 0.62 | 0.60 | 0.58 |
| CG Summary w/ Inline Citations (CG only) | 0.84 | 0.82 | 0.78 | 0.76 | 0.71 |
| PRISM-FC (PR + CG) | **0.86** | **0.84** | **0.83** | **0.81** | **0.76** |

### Table 3. Evidence category alignment (structured explainability)  
| Method | Category Alignment F1 (FR) | Category Alignment F1 (DE) |
|---|---:|---:|
| CG only | 0.67 | 0.65 |
| PRISM-FC (PR + CG) | **0.72** | **0.70** |
| PRISM-FC + Fairness/Toxicity Gate (FT) | 0.72 | 0.70 |

### Table 4. LLM-judge reliability under reference conflict (swapped-reference test)  
| Judge Setting | Correlation w/ Human (Normal Ref) | Correlation w/ Human (Swapped Ref) | Drop |
|---|---:|---:|---:|
| Standard LLM-as-Judge | 0.74 | 0.41 | -0.33 |
| Reference-Adherence Prompting | 0.73 | 0.52 | -0.21 |
| Reference-Adherence + NE2S scoring | **0.76** | **0.60** | **-0.16** |

---

## Discussion  
**Position bias matters downstream.** Improving retrieval robustness to evidence position (Table 1) yields consistent gains in citation precision and grounding-based content scores (Table 2). This supports PosIR’s central claim that position sensitivity is pervasive and consequential (Zeng et al., 2026), and shows it directly affects fact-checking summarization faithfulness.

**Inline citations substantially improve verifiability.** Enforcing citation-grounded generation raises citation precision markedly relative to unconstrained abstractive summaries (Table 2), aligning with the motivation of sui-1 for compliance-sensitive summarization (Droste et al., 2026). In fact-checking, this is critical because users must audit why a verdict is reached.

**Structured evidence categories improve interpretability.** Leveraging the dataset’s evidence schema (Hüsünbeyi et al., 2026) and improving retrieval increases category alignment (Table 3), suggesting that better evidence localization helps the model map evidence into meaningful, human-interpretable buckets.

**Evaluation must defend against judge failures.** The swapped-reference diagnostic reproduces the core failure mode identified by Lee et al. (2026): judge reliability collapses when references conflict with parametric knowledge. Our grounding-first NE2S scoring and reference-adherence prompting reduce—but do not eliminate—this vulnerability (Table 4). This implies that future benchmarks should include **reference-conflict stress tests** by default for LLM-judge evaluation in fact-checking.

**Fairness gate affects style more than task metrics.** The toxicity/fairness invocation does not change category alignment (Table 3), consistent with an inference-time refinement that targets consistency and fairness rather than accuracy (Ren et al., 2026). A limitation is that fairness impacts are not fully captured by the current metrics; targeted disparity measures would be needed.

---

## Conclusion  
We presented PRISM-FC, an end-to-end multilingual fact-checking summarization approach that combines position-robust retrieval, citation-grounded long-form synthesis, and evaluation diagnostics for LLM-judge reference adherence. Across French and German fact-checking articles, position-robust retrieval improves evidence recall and downstream citation precision, while citation constraints improve verifiability and grounding-based scores. Finally, swapped-reference tests reveal substantial judge unreliability, motivating grounding-first scoring and explicit reference-adherence protocols.

---

## Contributions  
1. **Problem framing:** Defined multilingual fact-checking long-form summarization as a **citation-grounded** and **evidence-category-structured** generation task grounded in authentic ClaimReview-derived debunking articles (Hüsünbeyi et al., 2026).  
2. **Method:** Proposed **PRISM-FC**, integrating **position-robust retrieval** (inspired by PosIR; Zeng et al., 2026) with **inline-cited long-form synthesis** (motivated by sui-1; Droste et al., 2026) and an inference-time **fairness/toxicity invocation gate** (Ren et al., 2026).  
3. **Evaluation protocol:** Introduced a **swapped-reference judge stress test** for fact-checking summarization evaluation, adapting the reference-conflict insights of Lee et al. (2026), and applied a grounding-first “No Evidence, No Score” principle inspired by evidence-chain evaluation (Ren et al., 2026).  
4. **Empirical findings:** Demonstrated that reducing retrieval position bias improves grounding and citation faithfulness, and that judge reliability under reference conflict remains a critical open challenge.

